% ============================================================
% 第2章 化学反应速率与反应历程（对应课件：选必一2.1 化学反应速率）
% ============================================================
\chapter{化学反应速率与反应历程}

\begin{daoxue}
\textbf{知识主线：}怎样描述快慢（速率的表示、计算、测定）$\to$ 什么因素影响快慢（浓度、压强、温度、催化剂）$\to$ 为什么会影响（碰撞理论、活化能）$\to$ 反应是怎样一步步发生的（基元反应、反应历程、催化机理）$\to$ 定量描述（速率方程、速率常数）。

\textbf{高考考情：}广东卷选择题第15题连续考查“反应历程能量图”（2024年考中间体、决速步与催化剂，2026年考平行反应的能量图与浓度—时间图）；原理大题常考速率计算和实验数据分析（2026第19题(3)）；实验题考控制变量与速率测定（2024第17题(3)）。

\textbf{学习目标：}①熟练进行速率计算与比较；②能用碰撞理论解释各因素的影响；③会读能量—历程图和催化循环图，找出决速步、中间体和催化剂；④理解速率方程与速率常数。
\end{daoxue}

\section{化学反应速率的表示与测定}

\begin{gainian}[化学反应速率]
用单位时间内反应物浓度的减少或生成物浓度的增加来表示（均取正值）：
\[v(\mathrm{B})=\frac{\Delta c(\mathrm{B})}{\Delta t}=\frac{\Delta n(\mathrm{B})}{V\cdot\Delta t}\]
常用单位 mol·L$^{-1}$·s$^{-1}$ 或 mol·L$^{-1}$·min$^{-1}$（$1\ \mathrm{mol\cdot L^{-1}\cdot s^{-1}}=60\ \mathrm{mol\cdot L^{-1}\cdot min^{-1}}$）。
\end{gainian}

理解这个定义要抓住三点：
\begin{enumerate}
\item 这样算出的是一段时间内的\textbf{平均速率}，不是某一时刻的瞬时速率（瞬时速率是 $c$–$t$ 曲线上该点切线斜率的绝对值）。
\item \textbf{固体和纯液体}的浓度视为常数，不能用它们的浓度变化表示速率（可用单位时间内质量或物质的量的变化表示）。
\item 同一反应 $m\mathrm{A}+n\mathrm{B}=p\mathrm{C}+q\mathrm{D}$，用不同物质表示的速率\textbf{数值可能不同，但意义相同}，且
\[v(\mathrm{A}):v(\mathrm{B}):v(\mathrm{C}):v(\mathrm{D})=m:n:p:q,\quad\text{即}\ \frac{v(\mathrm{A})}{m}=\frac{v(\mathrm{B})}{n}=\frac{v(\mathrm{C})}{p}=\frac{v(\mathrm{D})}{q}\]
因为在同一时间、同一体积内，各物质浓度变化量之比必然等于计量数之比。所以表示速率时必须指明是哪种物质。
\end{enumerate}

\begin{kanwu}[课件 2.1 第4页]
课件称“纯液体表示速率可用 L/s 或 L/min，纯固体可用 g/s 或 g/min”。规范表述为：纯液体、纯固体的浓度视为常数，不能用浓度变化表示速率；需要时可用单位时间内物质的量或质量的变化表示其反应快慢。
\end{kanwu}

\begin{fangfa}[比较速率快慢的“归一法”]
①统一单位；②都换算成同一种物质表示的速率，或计算 $\dfrac{v(\mathrm{X})}{\text{X的计量数}}$；③比较数值大小。
\end{fangfa}

\begin{liti}[速率比较]
对于 \ce{3A(g) + B(g) = 2C(g) + 2D(g)}，在不同条件下测得：①$v(\mathrm{A})=0.6$ mol·L$^{-1}$·min$^{-1}$；②$v(\mathrm{B})=0.45$ mol·L$^{-1}$·min$^{-1}$；③$v(\mathrm{C})=0.015$ mol·L$^{-1}$·s$^{-1}$；④$v(\mathrm{D})=0.45$ mol·L$^{-1}$·min$^{-1}$。反应最快的是\tk。
\tcblower
\jie 统一为 mol·L$^{-1}$·min$^{-1}$ 并除以计量数：①$0.6/3=0.2$；②$0.45/1=0.45$；③$0.015\times60/2=0.45$；④$0.45/2=0.225$。②与③相同且最快。\\
\xianjing ③若忘记把“s”换成“min”，会算出 0.0075 而误判为最慢。
\end{liti}

\begin{liti}[“三段式”计算]
容积为 2 L 的密闭容器中，将 2 mol L 气体和 3 mol M 气体混合，发生 \ce{2L(g) + 3M(g) = $x$Q(g) + 3R(g)}，10 s 末生成 2.4 mol R，并测得 Q 的浓度为 0.4 mol·L$^{-1}$。求：(1) 10 s 末 L 的浓度；(2) 前 10 s 内 $v(\mathrm{M})$；(3) $x$。
\tcblower
\silu 列“起始—转化—某时刻”三段，\textbf{只有转化量之比等于计量数之比}。\\
\jie \begin{tabular}[t]{lcccc}
 & 2L & 3M & $x$Q & 3R \\
起始/mol & 2 & 3 & 0 & 0 \\
转化/mol & 1.6 & 2.4 & 0.8 & 2.4 \\
10 s 末/mol & 0.4 & 0.6 & 0.8 & 2.4
\end{tabular}\\
(1) $c(\mathrm{L})=0.4/2=0.2$ mol·L$^{-1}$；(2) $v(\mathrm{M})=\dfrac{2.4}{2\times10}=0.12$ mol·L$^{-1}$·s$^{-1}$；(3) $n(\mathrm{Q})=0.4\times2=0.8$ mol，$x:3=0.8:2.4$，$x=1$。
\end{liti}

\textbf{速率的测定：}原则上，任何一种\textbf{随反应进行而明显变化、又便于测量}的量都可以用来测速率。例如锌与稀硫酸反应，可测一定时间内收集的 \ce{H2} 体积、收集一定体积 \ce{H2} 所需的时间、溶液 pH 的变化、锌粒质量的减少等；有颜色变化的反应可用比色法；离子浓度变化明显的可用电导法；生成沉淀的反应可测出现浑浊所需的时间；也可以每隔一段时间取样，用\textbf{滴定法}测某物质的浓度。

\begin{liti}[用量气法测定速率]
用锌粒与 40 mL 1 mol·L$^{-1}$ 稀硫酸反应，用注射器收集 \ce{H2}，测得 0～10 s 内收集到 \ce{H2} 44.8 mL（已换算为标准状况，溶液体积变化忽略不计）。(1) 用 \ce{H2} 的体积表示该段时间的速率；(2) 用 $c(\ce{H+})$ 的变化表示的速率为多少 mol·L$^{-1}$·s$^{-1}$？(3) 用 \ce{Zn^2+} 表示呢？
\tcblower
\jie (1) $v(\ce{H2})=\dfrac{44.8\ \mathrm{mL}}{10\ \mathrm{s}}=4.48$ mL·s$^{-1}$。(2) $n(\ce{H2})=2\times10^{-3}$ mol，消耗 \ce{H+} $4\times10^{-3}$ mol，$\Delta c(\ce{H+})=\dfrac{4\times10^{-3}}{0.040}=0.1$ mol·L$^{-1}$，$v(\ce{H+})=0.01$ mol·L$^{-1}$·s$^{-1}$。(3) $v(\ce{Zn^2+})=0.005$ mol·L$^{-1}$·s$^{-1}$（为 $v(\ce{H+})$ 的一半）。\\
\dianbo 用不同的量表示速率，单位不同、数值不同，但描述的是同一反应的快慢。
\end{liti}

\section{影响化学反应速率的因素}

\subsection{内因与外因}
决定反应速率的主要因素是\textbf{反应物本身的性质}（内因），如钠与水剧烈反应、镁与冷水反应很慢。对同一反应，外界条件（外因）也会显著改变速率：

\begin{center}
\small
\begin{tabularx}{\linewidth}{p{1.8cm}Xp{4.6cm}}
\toprule
条件 & 规律（其他条件不变时） & 说明 \\
\midrule
浓度 & 增大反应物浓度，速率增大 & 固体、纯液体的量不影响速率 \\
压强 & 对有气体参加的反应，压缩体积增大压强，速率增大 & 压强通过改变\textbf{浓度}起作用 \\
温度 & 升高温度，速率增大（每升高 10 ℃，速率通常增大到原来的 2～4 倍） & 正、逆反应速率都增大；与吸热放热无关 \\
催化剂 & 使用（正）催化剂，速率显著增大 & 同等程度改变正、逆反应速率 \\
其他 & 增大固体表面积（粉碎）、形成原电池、光照、超声波、搅拌等 & 原电池使锌与酸反应加快 \\
\bottomrule
\end{tabularx}
\end{center}

\begin{benzhi}[压强为什么只能“通过浓度”起作用]
对气体反应，速率取决于单位体积内的反应物分子数（浓度）。“增大压强”如果没有使反应物浓度变化，速率就不变。三种情况：
\begin{enumerate}
\item \textbf{恒温压缩体积}：各气体浓度都增大，正、逆反应速率都增大。
\item \textbf{恒温恒容充入惰性气体}（不参与反应的 He、Ar、\ce{N2} 等）：总压增大，但各反应物浓度不变，\textbf{速率不变}。
\item \textbf{恒温恒压充入惰性气体}：为保持压强不变，容器体积增大，各反应物浓度减小，\textbf{速率减小}。
\end{enumerate}
\end{benzhi}

\begin{liti}[惰性气体与速率]
在带活塞的密闭容器中发生 \ce{3Fe(s) + 4H2O(g) ->[\text{高温}] Fe3O4(s) + 4H2(g)}。判断下列操作对反应速率的影响：(1) 增加铁片的质量（表面积不变）；(2) 容积不变，充入更多水蒸气；(3) 压强不变，充入 Ar；(4) 容积不变，充入 Ar。
\tcblower
\jie (1) 不变：固体浓度视为常数。(2) 增大：反应物 \ce{H2O(g)} 浓度增大。(3) 减小：体积增大，\ce{H2O(g)}、\ce{H2(g)} 浓度减小。(4) 不变：各物质浓度均不变。
\end{liti}

\subsection{控制变量法探究影响因素}
研究某一因素对速率的影响时，必须保证\textbf{其他条件完全相同}，只改变这一个因素。溶液实验中，常通过\textbf{补加蒸馏水使总体积相同}，从而只改变某一反应物的浓度。

常用实验：①\ce{Na2S2O3 + H2SO4 = Na2SO4 + SO2 ^ + S v + H2O}，以“出现浑浊所需时间”衡量速率，探究浓度、温度的影响；②\ce{2MnO4^- + 5H2C2O4 + 6H+ = 2Mn^2+ + 10CO2 ^ + 8H2O}，以“褪色时间”衡量速率，需保证 \ce{KMnO4} 少量、\ce{H2C2O4} 过量，否则不褪色；③\ce{2H2O2 = 2H2O + O2 ^}，比较 \ce{FeCl3}、\ce{CuSO4} 的催化效果（为排除阴离子干扰，可改用 \ce{Fe2(SO4)3} 与 \ce{CuSO4} 比较）。

\begin{benzhi}[速率—时间曲线为什么会“先增大后减小”]
锌与足量盐酸反应：开始一段时间速率逐渐增大，是因为反应放热使溶液温度升高（温度的影响占主导）；后来速率减小，是因为 $c(\ce{H+})$ 减小（浓度的影响占主导）。类似地，\ce{KMnO4} 与 \ce{H2C2O4} 反应开始很慢、后来突然加快，是因为生成的 \ce{Mn^2+} 对反应有催化作用（\textbf{自催化}）。分析这类曲线时要找出“两种相反的影响”，再判断哪种占主导。
\end{benzhi}

\begin{zhenti}{2024年广东卷第17题(3)②\quad 来源：仓库所附真题文档}
（节选改述）探究 \ce{Na2CO3} 溶液浓度对反应 \ce{CaSO4(s) + CO3^2-(aq) <=> CaCO3(s) + SO4^2-(aq)} 速率的影响：已知 \ce{Na2CO3} 溶液浓度为 $\frac{c_1V_1}{V_0}$ mol·L$^{-1}$（$V_0$ mL 该溶液滴定消耗 $c_1$ mol·L$^{-1}$ 盐酸 $V_1$ mL，终点时 \ce{CO3^2-} 转化为 \ce{HCO3^-}）。按下表配制总体积相同的溶液，分别加入 $m_1$ g \ce{CaSO4} 固体，反应 $t_1$ min 后过滤，取 $V_0$ mL 滤液，用同一盐酸滴定，记录消耗盐酸的体积。

\begin{center}\small
\begin{tabular}{ccccc}
\toprule
序号 & $V(\ce{Na2CO3})$/mL & $V(\ce{H2O})$/mL & $V$(滤液)/mL & $V_\text{消耗}$(盐酸)/mL \\
\midrule
a & 100.0 & 0 & $V_0$ & $2V_1/5$ \\
b & 80.0 & $x$ & $V_0$ & $3V_1/10$ \\
\bottomrule
\end{tabular}
\end{center}
则 $x=$\tk，测得的平均反应速率之比 $v_\text{a}:v_\text{b}=$\tk。

\textbf{解析：}总体积相同，$x=20.0$。速率之比等于相同时间内 \ce{CO3^2-} 浓度减少量之比，以 $V_0$ mL 溶液中“消耗盐酸量”表示 \ce{CO3^2-}：a 中起始相当于 $c_1V_1$，剩余 $\frac25c_1V_1$，减少 $\frac35c_1V_1$；b 中起始为 $0.8c_1V_1$，剩余 $\frac{3}{10}c_1V_1$，减少 $\frac12c_1V_1$。$v_\text{a}:v_\text{b}=\frac35:\frac12=6:5$。
\end{zhenti}

\section{碰撞理论与活化能}

\subsection{有效碰撞与活化分子}
气体分子间的碰撞极其频繁（浓度 1 mol·L$^{-1}$ 时每立方厘米每秒约 $10^{28}$ 次），若每次碰撞都发生反应，所有气体反应都将瞬间完成。事实并非如此，说明\textbf{只有少数碰撞能引发反应}。

\begin{gainian}[有效碰撞、活化分子与活化能]
能够发生化学反应的碰撞叫\textbf{有效碰撞}，需同时满足：①分子具有足够高的能量；②碰撞时取向合适。能量足够、能发生有效碰撞的分子叫\textbf{活化分子}。活化分子具有的平均能量与反应物分子平均能量之差叫\textbf{活化能}（$E_\text{a}$）。活化能越小，活化分子所占比例越大，反应越快。
\end{gainian}

\begin{center}
\begin{tikzpicture}[scale=0.85]
\begin{scope}
\draw[-Stealth] (0,0)--(5.4,0) node[below left]{\small 分子能量};
\draw[-Stealth] (0,0)--(0,2.8) node[left]{\small 分子比例};
\fill[mainblue!30] (3.025,0) -- plot[domain=5.5:9,samples=40] ({0.55*\x},{4.0*\x*\x*exp(-\x)}) -- (4.95,0) -- cycle;
\fill[warnred!30] (2.2,0) -- plot[domain=4:5.5,samples=30] ({0.55*\x},{4.0*\x*\x*exp(-\x)}) -- (3.025,0) -- cycle;
\draw[thick] plot[domain=0:9,samples=90] ({0.55*\x},{4.0*\x*\x*exp(-\x)});
\draw[dashed] (3.025,0)--(3.025,1.0) node[above]{\scriptsize $E_\text{a}$};
\draw[dashed] (2.2,0)--(2.2,1.55) node[above]{\scriptsize $E_\text{a}'$};
\node[font=\scriptsize,align=left,anchor=west] at (3.0,2.3) {\textcolor{mainblue}{蓝}：无催化剂时的活化分子\\\textcolor{warnred}{红}$+$\textcolor{mainblue}{蓝}：有催化剂时的活化分子};
\end{scope}
\begin{scope}[xshift=8.8cm]
\axisEP{5.2}{3.2}
\draw[thick] (0.3,1.4)--(0.9,1.4) .. controls (1.7,1.4) and (1.9,2.8) .. (2.5,2.8) .. controls (3.1,2.8) and (3.3,0.6) .. (4.1,0.6)--(4.8,0.6);
\draw[thick,dashed,warnred] (0.9,1.4) .. controls (1.7,1.4) and (1.9,2.1) .. (2.5,2.1) .. controls (3.1,2.1) and (3.3,0.6) .. (4.1,0.6);
\draw[<->] (1.0,1.4)--(1.0,2.8) node[midway,left]{\scriptsize $E_1$};
\draw[<->] (4.5,0.6)--(4.5,2.8) node[midway,right]{\scriptsize $E_2$};
\draw[dotted] (1.0,2.8)--(4.5,2.8);
\node[font=\scriptsize] at (2.5,3.05) {过渡态};
\node[font=\scriptsize,warnred] at (3.2,1.9) {催化剂};
\end{scope}
\end{tikzpicture}
\end{center}
右图中 $E_1$ 为正反应活化能，$E_2$ 为逆反应活化能，$\Delta H=E_1-E_2$。催化剂改变反应途径，使 $E_1$、$E_2$ 降低相同的数值，所以 $\Delta H$ 不变。

\subsection{用碰撞理论解释外因的影响}
反应速率由“单位体积内单位时间的有效碰撞次数”决定，而它正比于\textbf{单位体积内的活化分子数}$=$单位体积内分子总数 $\times$ 活化分子百分数。

\begin{center}
\small
\begin{tabular}{lcccc}
\toprule
条件改变 & 单位体积分子总数 & 活化分子百分数 & 单位体积活化分子数 & 速率 \\
\midrule
增大反应物浓度（或压缩气体体积） & 增多 & 不变 & 增多 & 增大 \\
升高温度 & 不变 & 增大 & 增多 & 增大 \\
使用催化剂 & 不变 & 增大 & 增多 & 增大 \\
\bottomrule
\end{tabular}
\end{center}

\begin{yicuo}
\begin{enumerate}
\item 浓度（压强）只增加活化分子\textbf{数目}，不改变活化分子\textbf{百分数}；温度、催化剂才改变百分数。
\item 升高温度既有分子能量升高（更多分子成为活化分子，主要原因），也有碰撞频率增大（次要原因）。
\item 升温对吸热、放热反应都加快速率，对正、逆反应都加快；只是加快的程度不同，平衡才会移动（第3章）。
\end{enumerate}
\end{yicuo}

\section{基元反应、反应历程与催化机理}

\subsection{基元反应与反应历程}
大多数化学反应并非“一步到位”，而是经过若干个简单步骤才完成。每一步都叫一个\textbf{基元反应}；这些基元反应的总和叫\textbf{反应历程}（反应机理）。例如 \ce{2HI = H2 + I2} 的历程是：
\[\text{①}\ \ce{2HI -> H2 + 2I*}\qquad \text{②}\ \ce{2I* -> I2}\]
\ce{I*} 是带单电子的原子，称为\textbf{自由基}；它在①中生成、在②中消耗，是\textbf{中间产物}。

\begin{gainian}[决速步]
多步反应的总速率由其中\textbf{最慢}的一步决定，这一步称为决速步（速控步）。在能量—历程图上，决速步是\textbf{自身能垒（该步过渡态能量 $-$ 该步起始物质能量）最大}的那一步。
\end{gainian}

\begin{tuozhan}[飞秒化学——“看见”过渡态]
过渡态存在的时间极短（约 $10^{-13}$ s 量级）。化学家利用飞秒（1 fs $=10^{-15}$ s）激光脉冲，像“高速摄影”一样观测反应过程中化学键的断裂与形成，直接“看到”了反应物经过渡态转化为产物的过程（泽维尔因此获 1999 年诺贝尔化学奖）。这为“基元反应”“过渡态”等概念提供了实验依据。
\end{tuozhan}

\begin{yicuo}
判断决速步看的是“\textbf{每一步自己的爬坡高度}”，而不是哪个峰的绝对位置最高。例如某步起点很低、峰也较高，其能垒可能最大；另一步峰虽最高，但起点也高，能垒未必最大。数据常以 eV（单个分子）给出，读图时逐步做减法。
\end{yicuo}

\begin{liti}[多步反应的能量图]
某反应历程的相对能量如下（单位 eV）：反应物 0 $\to$ 过渡态 TS1 $+0.82$ $\to$ 中间体 M1 $-0.35$ $\to$ TS2 $+0.60$ $\to$ M2 $-0.20$ $\to$ TS3 $+0.25$ $\to$ 生成物 $-0.90$。(1) 决速步是第几步？其能垒为多少？(2) 该总反应吸热还是放热？(3) 能量最低的中间体是哪个？
\tcblower
\jie (1) 三步的能垒：第一步 $0.82-0=0.82$ eV；第二步 $0.60-(-0.35)=0.95$ eV；第三步 $0.25-(-0.20)=0.45$ eV。决速步为第二步（\ce{M1 -> M2}），能垒 0.95 eV——虽然 TS1 的绝对能量最高，但它不是决速步。(2) 生成物能量 $-0.90$ eV 低于反应物，放热。(3) M1（$-0.35$ eV）。
\end{liti}

\subsection{催化剂的作用机理}
催化剂参与反应、改变反应历程，但反应前后其质量和化学性质不变。例如 \ce{Fe+} 催化 \ce{N2O} 与 CO 的反应：
\[\text{①}\ \ce{Fe+ + N2O -> FeO+ + N2}\qquad\text{②}\ \ce{FeO+ + CO -> Fe+ + CO2}\qquad\text{总：}\ \ce{N2O + CO ->[\ce{Fe+}] N2 + CO2}\]

\begin{fangfa}[区分催化剂与中间产物]
\textbf{催化剂}：先参与反应（出现在前面某步的反应物中），后又再生（出现在后面某步的生成物中），总方程式中不出现，如上例的 \ce{Fe+}。\textbf{中间产物}：先生成、后消耗，总方程式中也不出现，如上例的 \ce{FeO+}。在“循环图”中，从外部进入循环的是\textbf{反应物}，离开循环的是\textbf{生成物}，循环圈上的物种中，作为循环起点并最终回到原状的是\textbf{催化剂}，其余是中间体。
\end{fangfa}

\begin{liti}[读催化循环图]
某金属催化剂 M 催化 \ce{X + Y -> Z} 的循环机理为：①\ce{M + X -> M-X}；②\ce{M-X + Y -> M-XY}；③\ce{M-XY -> M + Z}。
\begin{center}
\begin{tikzpicture}[font=\small,scale=0.9]
\node (M) at (90:1.5) {M}; \node (MX) at (-30:1.5) {M—X}; \node (MXY) at (210:1.5) {M—XY};
\draw[-Stealth,thick,mainblue] (M) to[bend left=30] node[right=3pt]{\scriptsize ① $+$X} (MX);
\draw[-Stealth,thick,mainblue] (MX) to[bend left=30] node[below]{\scriptsize ② $+$Y} (MXY);
\draw[-Stealth,thick,mainblue] (MXY) to[bend left=30] node[left=3pt]{\scriptsize ③ $-$Z} (M);
\end{tikzpicture}
\end{center}
(1) 指出反应物、生成物、催化剂和中间体；(2) 若第②步最慢，增大 $c(\mathrm{X})$ 对总速率影响大还是增大 $c(\mathrm{Y})$ 影响大？
\tcblower
\jie (1) 从外部进入循环的 X、Y 是反应物，离开循环的 Z 是生成物；M 先消耗后再生，是催化剂；M—X、M—XY 先生成后消耗，是中间体。(2) 总速率由第②步决定，其速率与 $c(\ce{M-X})$、$c(\mathrm{Y})$ 有关；直接增大 $c(\mathrm{Y})$ 能加快决速步，影响更大（增大 $c(\mathrm{X})$ 只有在第①步建立的 M—X 浓度随之增大时才有作用）。
\end{liti}

\begin{benzhi}[催化剂“能做什么、不能做什么”]
\textbf{能}：降低活化能、加快正逆反应速率（同等倍数）、缩短达到平衡的时间；有\textbf{选择性}——在多个可能的反应中只加快某一个，从而提高目标产物的选择性。\textbf{不能}：改变 $\Delta H$、改变平衡常数和平衡转化率、使不能发生（非自发）的反应发生。此外，催化剂需要合适的温度（活性温度），温度过高可能失活；某些杂质会使催化剂“中毒”，所以工业原料气要净化。
\end{benzhi}

\begin{zhenti}{2024年广东卷第15题\quad 来源：仓库所附真题文档}
对反应 \ce{S(g) <=> T(g)}（I 为中间产物），相同条件下：①加入催化剂，反应达到平衡所需时间大幅缩短；②提高反应温度，$c_\text{平}(\mathrm{S})/c_\text{平}(\mathrm{T})$ 增大，$c_\text{平}(\mathrm{S})/c_\text{平}(\mathrm{I})$ 减小。基于以上事实，可能的反应历程示意图（实线为无催化剂，虚线为有催化剂）为
\begin{center}
\begin{tikzpicture}[xscale=0.42,yscale=0.42,font=\scriptsize]
\foreach \lab/\xs/\Ilev/\cat in {A/0/1.6/1, B/7.5/1.6/2, C/15/0.4/1, D/22.5/0.4/2}{
\begin{scope}[xshift=\xs cm]
\draw[-Stealth] (0,0)--(6.2,0); \draw[-Stealth] (0,0)--(0,5.2);
\node[left] at (-0.2,4.6) {\lab.};
\draw[thick] (0.3,1.1)--(0.6,1.1) .. controls (1.2,1.1) and (1.3,4.6) .. (1.9,4.6) .. controls (2.5,4.6) and (2.7,\Ilev) .. (3.1,\Ilev) .. controls (3.5,\Ilev) and (3.7,3.6) .. (4.1,3.6) .. controls (4.5,3.6) and (4.8,0.3) .. (5.3,0.3)--(5.7,0.3);
\ifnum\cat=1 \draw[dashed,thick] (0.6,1.1) .. controls (1.0,1.1) and (1.1,2.9) .. (1.5,2.9) .. controls (1.8,2.9) and (2.0,2.5) .. (2.2,2.5) .. controls (2.4,2.5) and (2.5,3.4) .. (2.7,3.4) .. controls (2.9,3.4) and (2.95,\Ilev) .. (3.1,\Ilev);
\else \draw[dashed,thick] (3.1,\Ilev) .. controls (3.4,\Ilev) and (3.5,2.6) .. (3.7,2.6) .. controls (3.85,2.6) and (3.9,2.2) .. (4.05,2.2) .. controls (4.2,2.2) and (4.3,2.7) .. (4.45,2.7) .. controls (4.7,2.7) and (4.85,0.3) .. (5.3,0.3);
\fi
\node[below] at (0.45,1.1) {S}; \node[below] at (3.1,\Ilev) {I}; \node[below] at (5.5,0.3) {T};
\end{scope}}
\end{tikzpicture}
\end{center}
\textbf{解析：}②中升温 $c(\mathrm{S})/c(\mathrm{T})$ 增大，说明 \ce{S <=> T} 平衡逆移，正反应放热，T 的能量低于 S；升温 $c(\mathrm{S})/c(\mathrm{I})$ 减小，说明 \ce{S <=> I} 正向吸热，I 的能量高于 S，排除 C、D（I 低于 S）。①中加入催化剂使达到平衡的时间\textbf{大幅}缩短，说明催化剂降低的是决速步（能垒大的第一步）的活化能，A 正确；B 中催化剂只降低了能垒较小的第二步，对总速率影响不大。答案：A。
\end{zhenti}

\begin{zhenti}{2026年广东卷第15题\quad 来源：仓库所附真题文档（网络收集版）}
物质 A(g) 可同时转化为它的同分异构体 B(g) 和 C(g)。图a（能量—反应过程）显示：A 生成 C 的活化能较小，C 的能量略低于 A，B 的能量最低；图b 为某温度下恒容密闭容器中各组分浓度随时间的变化：$t_1$ 时 B、C 浓度相等（均为 $x$），A 的浓度为 $y$；$t_2$ 时 A、B 浓度相等（均为 $z$）；C 的浓度先增大后减小。则
\xxb{$t_1$ 时，$m(\mathrm{A})=2m(\mathrm{B})$}{$2z>x+y$}{体系中，反应 \ce{A(g) <=> C(g)} 先达到平衡}{体系平衡后，升高温度，B 的体积分数增加}
\textbf{解析：}设 A 起始浓度为 $c_0$，三者互为同分异构体，由原子守恒，任意时刻 $c(\mathrm{A})+c(\mathrm{B})+c(\mathrm{C})=c_0$。$t_1$：$y+2x=c_0$，A 与 B 的质量关系无法确定，A 错。$t_2$：设 $c(\mathrm{C})=a$，$2z+a=c_0=2x+y$，得 $2z=x+y+(x-a)$；C 先增后减，且 $t_2$ 时 C 已在减小，$x>a$，故 $2z>x+y$，B 正确。两个反应共用反应物 A，A 的浓度不变时两个平衡同时建立，不存在谁“先达到平衡”，C 错。B 的能量低于 A，\ce{A -> B} 放热，升温平衡逆移，B 的体积分数减小，D 错。答案：B。

\dianbo C 生成快（活化能小）但最终不如 B 多（B 更稳定）：这体现了\textbf{动力学（速率）控制与热力学（平衡）控制}的区别——短时间内产物以“生成快”的为主，足够长时间后以“更稳定”的为主。
\end{zhenti}

\section{速率方程与速率常数}

\begin{gainian}[速率方程]
一定温度下，反应 $a\mathrm{A}+b\mathrm{B}\to$ 产物 的瞬时速率常可表示为
\[v=k\cdot c^{\,m}(\mathrm{A})\cdot c^{\,n}(\mathrm{B})\]
$k$ 为\textbf{速率常数}，与温度、催化剂、固体表面性质等有关，\textbf{与浓度无关}；$m$、$n$ 由实验测定，对\textbf{基元反应}才等于计量数（即 $m=a$，$n=b$）。$k$ 越大，反应越快。
\end{gainian}

\textbf{确定速率方程的方法：}控制变量——保持其他物质浓度不变，某物质浓度变为原来的 2 倍时，若速率变为 2 倍，则其指数为 1；变为 4 倍，指数为 2；不变，指数为 0。

\begin{benzhi}[温度如何影响速率常数——阿伦尼乌斯公式]
$k=A\mathrm{e}^{-E_\text{a}/(RT)}$，取对数得 $\ln k=\ln A-\dfrac{E_\text{a}}{R}\cdot\dfrac1T$。以 $\ln k$（或 $R\ln k$）对 $\dfrac1T$ 作图得直线，\textbf{斜率 $=-\dfrac{E_\text{a}}{R}$}（或 $-E_\text{a}$）。由此可知：①升高温度，$k$ 增大；②$E_\text{a}$ 越大，直线越陡，速率对温度越敏感；③催化剂降低 $E_\text{a}$，直线变平缓。高考中常给出 $R\ln k$–$\frac1T$ 图，要求由斜率求活化能或比较催化剂优劣，只需会读图，不需记忆公式。
\end{benzhi}

\begin{zhenti}{2026年广东卷第19题(3)\quad 来源：仓库所附真题文档（网络收集版）}
\ce{H2S} 的氧化可在卟啉铁配合物 [\ce{FeL(H2O)}]（用 Z 表示）催化下进行，其中一步为 \ce{[FeL(H2O)] + HS- <=> [FeL(HS)]- + H2O}。为研究不同因素对该反应的影响，进行 4 组实验：Ⅰ～Ⅲ 组仅初始 $c_0(\ce{HS-})$ 不同（$c_\text{Ⅰ}<c_\text{Ⅱ}<c_\text{Ⅲ}$）；Ⅳ 与 Ⅱ 组仅温度不同。$0\sim t$ 时间内 Z 的平均反应速率 $-\frac{\Delta c(\mathrm{Z})}{\Delta t}$（单位 $10^{-6}$ mol·L$^{-1}$·s$^{-1}$）：Ⅰ 为 1.0，Ⅱ 为 2.0，Ⅲ 为 $y$，Ⅳ 为 1.0。\\
①下列说法正确的有\tk。A．$0\sim t$ 内反应速率：Ⅰ$>$Ⅱ$>$Ⅲ\quad B．Z 的平衡转化率：Ⅰ$<$Ⅲ\quad C．实验温度：Ⅱ$>$Ⅳ\quad D．平衡时 Z 的浓度：Ⅱ 和 Ⅳ 相等\\
②Ⅰ、Ⅱ 和 Ⅲ 的数据点分布在一条直线上，且 $(c_\text{Ⅲ}-c_\text{Ⅱ})/(c_\text{Ⅱ}-c_\text{Ⅰ})=1.2$，则 $y=$\tk。

\textbf{解析：}①随 $c_0(\ce{HS-})$ 增大，Z 的反应速率增大，A 错；另一反应物 \ce{HS-} 浓度越大，Z 的平衡转化率越大，B 对；Ⅱ、Ⅳ 浓度相同而 Ⅱ 速率大，温度更高，C 对；温度不同，平衡常数不同，平衡时 Z 的浓度不相等，D 错。选 BC。②三点共线，斜率相同：$\dfrac{y-2.0}{c_\text{Ⅲ}-c_\text{Ⅱ}}=\dfrac{2.0-1.0}{c_\text{Ⅱ}-c_\text{Ⅰ}}$，$y-2.0=1.2$，$y=3.2$。
\end{zhenti}

\textbf{速率常数与平衡常数：}对基元可逆反应 \ce{$a$A + $b$B <=> $c$C + $d$D}，$\vz=\kz c^a(\mathrm{A})c^b(\mathrm{B})$，$\vn=\kn c^c(\mathrm{C})c^d(\mathrm{D})$。平衡时 $\vz=\vn$，得 $K=\dfrac{\kz}{\kn}$。升高温度，$\kz$、$\kn$ 都增大；若正反应放热，则 $\kn$ 增大得更多，$K$ 减小（详见第3章）。

% ------------------------------------------------------------
\zhangmo{本章小结}
\subsection*{一、知识总结}
\begin{center}
\small
\begin{tabularx}{\linewidth}{p{2.4cm}X}
\toprule
速率表示 & $v=\Delta c/\Delta t$（平均速率）；速率之比$=$计量数之比；固体、纯液体不用浓度表示 \\
影响因素 & 内因决定；浓度、压强（通过浓度）、温度、催化剂、表面积等外因 \\
碰撞理论 & 浓度改变活化分子数；温度、催化剂改变活化分子百分数 \\
反应历程 & 基元反应；决速步$=$自身能垒最大的一步；催化剂先消耗后再生，中间体先生成后消耗 \\
速率方程 & $v=kc^m(\mathrm{A})c^n(\mathrm{B})$，指数实验测定；$k$ 只与温度、催化剂等有关；基元反应 $K=\kz/\kn$ \\
\bottomrule
\end{tabularx}
\end{center}

\subsection*{二、方法提升}
\begin{enumerate}
\item \textbf{速率计算}：一律先列三段式，只用“转化量”与计量数成比例。
\item \textbf{惰性气体问题}：先判断容器是恒容还是恒压，再看反应物浓度是否改变。
\item \textbf{历程图三问}：吸放热看起终点；决速步看每步能垒；催化剂、中间体看“先进后出”还是“先出后进”。
\item \textbf{数据分析题}：用控制变量找“只有一个条件不同”的两组实验对比；涉及多组数据时，注意线性关系、倍数关系。
\end{enumerate}

\subsection*{三、易错点归纳}
\begin{enumerate}
\item 误用固体或纯液体的浓度变化计算速率；比较速率时忘记统一单位。
\item 认为“增大压强速率一定增大”，忽略恒容充入惰性气体的情况。
\item 认为升温只加快吸热反应；认为催化剂只加快正反应。
\item 把能量最高的过渡态所在步骤当作决速步；把中间体当作催化剂。
\item 认为速率方程中的指数一定等于方程式计量数（只有基元反应才是）。
\end{enumerate}

% ------------------------------------------------------------
\zhangmo{分层练习}
\zu{A组\quad 基础巩固}
\begin{xiti}
\item 下列关于化学反应速率的说法错误的是\xz
\xxa{化学反应速率是衡量化学反应进行快慢的物理量}{决定化学反应速率的主要因素是反应物本身的性质}{某反应在一段时间内的速率为 0.5 mol·L$^{-1}$·s$^{-1}$，是指该段时间内反应物和生成物的浓度变化都为 0.5 mol·L$^{-1}$}{可以用单位时间内固体质量的变化来表示有固体参加的反应的快慢}

\item 合成氨 \ce{N2 + 3H2 <=> 2NH3} 在不同条件下测得：①$v(\ce{N2})=0.6$ mol·L$^{-1}$·min$^{-1}$；②$v(\ce{H2})=1.2$ mol·L$^{-1}$·min$^{-1}$；③$v(\ce{NH3})=0.03$ mol·L$^{-1}$·s$^{-1}$。反应速率由大到小的顺序为\tk。

\item 一定温度下，在 2 L 密闭容器中充入 1 mol \ce{H2(g)} 和 1 mol \ce{I2(g)}，发生 \ce{H2(g) + I2(g) <=> 2HI(g)} $\Delta H<0$。下列说法正确的是\xz
\xxa{保持容积不变，再充入 1 mol \ce{H2(g)}，反应速率加快}{保持容积不变，充入 1 mol \ce{N2}（不参与反应），反应速率加快}{保持压强不变，充入 1 mol \ce{N2}（不参与反应），反应速率不变}{保持压强不变，再充入 1 mol \ce{H2(g)} 和 1 mol \ce{I2(g)}，反应速率加快}

\item 下列有关有效碰撞理论的说法正确的是\xz
\xxa{增大反应物浓度，可增大活化分子百分数}{升高温度，活化分子百分数增大，有效碰撞频率增大}{催化剂能降低反应的活化能，因此能使 $\Delta H$ 减小}{活化分子之间的碰撞一定是有效碰撞}

\item 下列事实能说明影响化学反应速率的决定性因素是反应物本身性质的是\xz
\xxa{大小、形状相同的镁片、铁片分别与同浓度盐酸反应，前者产生 \ce{H2} 更快}{Cu 与浓硝酸反应比与稀硝酸反应快}{\ce{N2} 与 \ce{O2} 在常温常压下不反应，放电时可反应}{Cu 能与浓硫酸反应，而不与稀硫酸反应}

\item 下列关于催化剂的说法正确的是\xz
\xxa{催化剂不参加化学反应，反应前后质量不变}{催化剂可以提高反应物的平衡转化率}{同一可逆反应中，催化剂对正、逆反应速率的改变程度相同}{任何温度下催化剂的催化效果都相同}

\item 已知温度每升高 10 ℃，某反应速率增大为原来的 2 倍。若将温度由 20 ℃ 升高到 50 ℃，反应速率增大为原来的\tk{}倍；若原来反应需 16 min 完成，50 ℃ 时约需\tk min。
\end{xiti}

\zu{B组\quad 高考典型}
\begin{xiti}
\item \laiyuan{原创}\ 用 0.10 mol·L$^{-1}$ \ce{Na2S2O3} 溶液与 0.10 mol·L$^{-1}$ \ce{H2SO4} 溶液探究影响速率的因素，记录出现浑浊所需时间：
\begin{center}\small
\begin{tabular}{cccccc}
\toprule
实验 & 温度/℃ & $V(\ce{Na2S2O3})$/mL & $V(\ce{H2SO4})$/mL & $V(\ce{H2O})$/mL & 出现浑浊时间/s \\
\midrule
① & 20 & 10.0 & 10.0 & 0 & $t_1$ \\
② & 20 & 5.0 & 10.0 & $a$ & $t_2$ \\
③ & 40 & 10.0 & 10.0 & $b$ & $t_3$ \\
\bottomrule
\end{tabular}
\end{center}
(1) 写出反应的离子方程式。(2) $a=$\tk，$b=$\tk。(3) 实验①②探究的是\tk 对速率的影响，预测 $t_1$\tk$t_2$（填“$>$”“$<$”或“$=$”）。(4) 实验中能否用“生成 \ce{SO2} 气体的体积”来衡量速率？说明理由。

\item \laiyuan{原创}\ \ce{Fe+} 可催化 \ce{N2O(g) + CO(g) = N2(g) + CO2(g)}，反应分两步：①\ce{Fe+ + N2O -> FeO+ + N2}；②\ce{FeO+ + CO -> Fe+ + CO2}。其能量变化示意图中，第①步活化能为 $E_1$，第②步活化能为 $E_2$，且 $E_1>E_2$，两步均放热。下列说法错误的是\xz
\xxa{\ce{Fe+} 是催化剂，\ce{FeO+} 是中间产物}{总反应速率主要由第①步决定}{使用 \ce{Fe+} 后总反应的 $\Delta H$ 减小}{第①步反应的速率比第②步慢}

\item 已知 \ce{2NO(g) + 2H2(g) = N2(g) + 2H2O(g)}，800 ℃ 时测得如下数据：
\begin{center}\small
\begin{tabular}{cccc}
\toprule
实验 & $c(\ce{NO})$/(mol·L$^{-1}$) & $c(\ce{H2})$/(mol·L$^{-1}$) & 初始速率/(mol·L$^{-1}$·s$^{-1}$) \\
\midrule
1 & 0.10 & 0.10 & $2.0\times10^{-3}$ \\
2 & 0.20 & 0.10 & $8.0\times10^{-3}$ \\
3 & 0.10 & 0.20 & $4.0\times10^{-3}$ \\
\bottomrule
\end{tabular}
\end{center}
(1) 写出速率方程，并求速率常数 $k$（注明单位）。(2) 求 $c(\ce{NO})=0.30$ mol·L$^{-1}$、$c(\ce{H2})=0.05$ mol·L$^{-1}$ 时的速率。(3) 有人提出该反应的历程：Ⅰ．\ce{2NO <=> N2O2}（快）；Ⅱ．\ce{N2O2 + H2 -> N2O + H2O}（慢）；Ⅲ．\ce{N2O + H2 -> N2 + H2O}（快）。指出决速步和中间产物。

\item \laiyuan{原创}\ 某反应在两种催化剂 a、b 作用下，$\ln k$ 与 $\frac1T$ 的关系为直线：催化剂 a 的直线过点 $(2.0\times10^{-3},\,2.0)$ 和 $(3.0\times10^{-3},\,-4.0)$；催化剂 b 的直线过点 $(2.0\times10^{-3},\,3.0)$ 和 $(3.0\times10^{-3},\,0)$（$\frac1T$ 单位 K$^{-1}$）。已知 $\ln k=\ln A-\frac{E_\text{a}}{R}\cdot\frac1T$，$R=8.314$ J·mol$^{-1}$·K$^{-1}$。(1) 分别求两种催化剂下的活化能（kJ·mol$^{-1}$，保留三位有效数字）。(2) 哪种催化剂效果更好？哪种情况下反应速率对温度更敏感？

\item \laiyuan{原创}\ 用酸性 \ce{KMnO4} 溶液与 \ce{H2C2O4} 溶液反应探究影响速率的因素（\ce{KMnO4} 均少量，溶液总体积相同）：
\begin{center}\small
\begin{tabular}{ccccc}
\toprule
实验 & 温度/℃ & 加入的其他物质 & $c(\ce{H2C2O4})$ & 褪色时间/s \\
\midrule
① & 25 & 无 & 0.10 mol·L$^{-1}$ & 40 \\
② & 25 & 少量 \ce{MnSO4} 固体 & 0.10 mol·L$^{-1}$ & 4 \\
③ & 25 & 无 & 0.20 mol·L$^{-1}$ & 25 \\
④ & 50 & 无 & 0.10 mol·L$^{-1}$ & 10 \\
\bottomrule
\end{tabular}
\end{center}
(1) 写出反应的离子方程式。(2) 对比实验\tk{}可探究浓度的影响，对比实验\tk{}可探究温度的影响。(3) 实验②说明了什么？据此解释实验①中“开始褪色很慢，一段时间后突然加快”的现象。

\item \laiyuan{原创}\ 某反应在催化剂 A、B 作用下的历程如下（相对能量，eV）：催化剂 A：反应物 0 $\to$ TS1 $+1.10$ $\to$ 中间体 $+0.20$ $\to$ TS2 $+0.85$ $\to$ 产物 $-0.50$；催化剂 B：反应物 0 $\to$ TS1$'$ $+0.70$ $\to$ 中间体 $-0.40$ $\to$ TS2$'$ $+0.60$ $\to$ 产物 $-0.50$。(1) 分别指出两种催化剂下的决速步及其能垒。(2) 哪种催化剂效果更好？(3) 两种催化剂下反应的 $\Delta H$ 是否相同？
\end{xiti}

\zu{C组\quad 冲刺高分}
\begin{xiti}
\item \laiyuan{原创}\ 某温度下，向恒容密闭容器中充入 A(g)，同时发生两个反应：\ce{A(g) <=> B(g)}（反应Ⅰ，活化能大，$\Delta H_1<0$）和 \ce{A(g) <=> C(g)}（反应Ⅱ，活化能小，$\Delta H_2<0$，且 $|\Delta H_2|<|\Delta H_1|$）。
(1) 反应开始一段时间内，B、C 中哪一种产物的浓度增大得快？为什么？
(2) 足够长时间后，平衡体系中 B、C 哪一种含量高？为什么？
(3) 若要在较短时间内获得尽可能多的 C，应如何控制反应时间？若要获得 B，应如何操作？
(4) 若选用一种只加快反应Ⅰ的催化剂，平衡时 B 的体积分数是否改变？说明理由。

\item \laiyuan{原创}\ 在恒容密闭容器中发生基元反应 \ce{2NO2(g) <=> N2O4(g)} $\Delta H<0$，$\vz=\kz c^2(\ce{NO2})$，$\vn=\kn c(\ce{N2O4})$。$T_1$ 时，$\kz=2\kn$（浓度单位 mol·L$^{-1}$）。
(1) $T_1$ 时该反应的平衡常数 $K=$\tk。
(2) $T_1$ 时，若某时刻 $c(\ce{NO2})=0.5$ mol·L$^{-1}$，$c(\ce{N2O4})=1.0$ mol·L$^{-1}$，判断此时 $\vz$ 与 $\vn$ 的大小。
(3) 升高温度至 $T_2$，$\kz$、$\kn$ 如何变化？$\kz/\kn$ 增大还是减小？

\item \laiyuan{原创}\ 反应 \ce{2NO(g) + O2(g) = 2NO2(g)} 的速率方程为 $v=kc^2(\ce{NO})c(\ce{O2})$。有人提出两步机理：Ⅰ．\ce{2NO <=> N2O2}（快速平衡，$\Delta H<0$）；Ⅱ．\ce{N2O2 + O2 -> 2NO2}（慢）。(1) 指出决速步和中间产物。(2) 用第Ⅰ步的平衡常数 $K_1$ 和第Ⅱ步的速率常数 $k_2$ 表示总反应的 $k$。(3) 实验发现：升高温度，该反应的速率反而略有减小。结合机理解释。
\end{xiti}

% ------------------------------------------------------------
\zhangmo{练习解析}
\begin{jiexilist}
\item \ans{C}速率是单位时间内浓度的变化量，单位是 mol·L$^{-1}$·s$^{-1}$；而且用不同物质表示时数值不一定相同，C 错。D 对：可用单位时间内固体质量（或物质的量）的变化表示快慢，只是不能用其“浓度”表示。
\item \ans{③$>$①$>$②}都换算为 $v(\ce{N2})$（mol·L$^{-1}$·min$^{-1}$）：①0.6；②$1.2/3=0.4$；③$0.03\times60/2=0.9$。注意③的单位是“s”，必须先乘 60 换算为“min”。
\item \ans{A}A：反应物 \ce{H2} 浓度增大，速率加快。B：恒容充 \ce{N2}，各物质浓度不变，速率不变。C：恒压充 \ce{N2}，体积增大，浓度减小，速率减小。D：恒压下按原比例再充入 \ce{H2}、\ce{I2}，体积同比例增大，各物质浓度与原来相同，速率不变。
\item \ans{B}A 错：浓度增大的是单位体积活化分子数，百分数不变。C 错：催化剂不改变 $\Delta H$。D 错：活化分子还需取向合适才能发生有效碰撞。
\item \ans{A}A 中两种金属本身性质不同，其他条件相同，体现内因。B 是浓度的影响；C 是放电（外界能量）的影响；D 是浓度不同导致性质不同（浓硫酸有强氧化性），不属于“不同反应物本身性质”的比较。
\item \ans{C}A 错：催化剂参与反应，只是反应前后质量和化学性质不变。B 错：催化剂不能使平衡移动。D 错：催化剂有活性温度范围。
\item 8；2。升高 30 ℃，即 3 个 10 ℃，速率变为 $2^3=8$ 倍，时间变为原来的 $\frac18$。
\item (1) \ce{S2O3^2- + 2H+ = S v + SO2 ^ + H2O}。(2) $a=5.0$，$b=0$：①总体积 20.0 mL，②中补水 5.0 mL 使总体积相同；③只改变温度，不加水。(3) \ce{Na2S2O3} 浓度；$t_1<t_2$（①浓度大，速率快，出现浑浊所需时间短）。(4) 不宜。\ce{SO2} 易溶于水，生成的 \ce{SO2} 有相当一部分溶于溶液，测得的气体体积不能准确反映反应进程。
\item \ans{C}\ce{Fe+} 先消耗后再生，是催化剂；\ce{FeO+} 先生成后消耗，是中间产物，A 对。$E_1>E_2$，第①步慢，决定总速率，B、D 对。催化剂不改变 $\Delta H$，C 错。
\item (1) 由实验1、2：$c(\ce{NO})$ 加倍，速率变为 4 倍，NO 的指数为 2；由实验1、3：$c(\ce{H2})$ 加倍，速率加倍，\ce{H2} 的指数为 1。$v=kc^2(\ce{NO})c(\ce{H2})$；$k=\dfrac{2.0\times10^{-3}}{0.10^2\times0.10}=2.0$ L$^2$·mol$^{-2}$·s$^{-1}$。(2) $v=2.0\times0.30^2\times0.05=9.0\times10^{-3}$ mol·L$^{-1}$·s$^{-1}$。(3) 决速步为Ⅱ（慢反应）；中间产物为 \ce{N2O2} 和 \ce{N2O}。（说明：第Ⅰ步为快速平衡，\ce{N2O2} 浓度正比于 $c^2(\ce{NO})$，所以慢步Ⅱ的速率 $\propto c^2(\ce{NO})c(\ce{H2})$，与实验测得的速率方程一致，这正是用实验检验机理的思路。）
\item (1) 斜率 $=-\dfrac{E_\text{a}}{R}$。a：斜率 $=\dfrac{-4.0-2.0}{1.0\times10^{-3}}=-6.0\times10^3$ K，$E_\text{a}=6.0\times10^3\times8.314$ J·mol$^{-1}\approx49.9$ kJ·mol$^{-1}$；b：斜率 $=\dfrac{0-3.0}{1.0\times10^{-3}}=-3.0\times10^3$ K，$E_\text{a}\approx24.9$ kJ·mol$^{-1}$。(2) b 的活化能小，且在所给温度范围内 $\ln k$ 都更大，催化效果更好；a 的直线更陡（$E_\text{a}$ 更大），使用 a 时速率对温度更敏感。
\item (1) \ce{2MnO4^- + 5H2C2O4 + 6H+ = 2Mn^2+ + 10CO2 ^ + 8H2O}。(2) ①③；①④。(3) \ce{Mn^2+} 对该反应有催化作用；实验①开始时 \ce{Mn^2+} 很少，反应很慢，随着反应生成 \ce{Mn^2+}（自催化），速率突然加快。
\item (1) A：第一步能垒 1.10 eV，第二步 $0.85-0.20=0.65$ eV，决速步为第一步（1.10 eV）；B：第一步 0.70 eV，第二步 $0.60-(-0.40)=1.00$ eV，决速步为第二步（1.00 eV）。(2) B 的最大能垒（1.00 eV）小于 A（1.10 eV），B 效果更好。(3) 相同，都为 $-0.50$ eV（每个分子），$\Delta H$ 与催化剂无关。注意：B 的 TS2$'$ 绝对能量虽低于 A 的 TS1，但其决速步是第二步——必须逐步计算能垒。
\item (1) C。反应Ⅱ活化能小，速率快，开始阶段 C 生成得快。(2) B。两反应都放热，$|\Delta H_1|>|\Delta H_2|$，B 的能量最低、最稳定，平衡（足够长时间后）B 的含量高；此时生成的 C 会经 A 逐步转化为 B。(3) 获得 C：控制较短的反应时间，在 C 浓度达到最大值附近时及时终止反应（分离产物）；获得 B：延长反应时间使体系接近平衡（也可选用能加快反应Ⅰ的催化剂）。(4) 不改变。催化剂不改变平衡常数，平衡状态只由温度和起始组成决定，B 的平衡体积分数不变；催化剂只缩短 B 达到平衡含量所需的时间。
\item (1) 2。平衡时 $\kz c^2(\ce{NO2})=\kn c(\ce{N2O4})$，$K=\dfrac{c(\ce{N2O4})}{c^2(\ce{NO2})}=\dfrac{\kz}{\kn}=2$。(2) $Q=\dfrac{1.0}{0.5^2}=4>K=2$，反应逆向进行，$\vz<\vn$（也可直接算：$\vz=2\kn\times0.25=0.5\kn$，$\vn=\kn\times1.0$，$\vz<\vn$）。(3) 升温，$\kz$、$\kn$ 都增大；正反应放热，升温平衡逆移，$K$ 减小，故 $\kz/\kn$ 减小（$\kn$ 增大的倍数更大）。
\item (1) 第Ⅱ步为决速步；中间产物为 \ce{N2O2}。(2) 第Ⅰ步快速平衡：$c(\ce{N2O2})=K_1c^2(\ce{NO})$；总速率由慢步决定：$v=k_2c(\ce{N2O2})c(\ce{O2})=k_2K_1c^2(\ce{NO})c(\ce{O2})$，所以 $k=k_2K_1$。(3) 升高温度，$k_2$ 增大；但第Ⅰ步放热，升温使 $K_1$ 减小，$c(\ce{N2O2})$ 降低。若 $K_1$ 减小的程度超过 $k_2$ 增大的程度，$k=k_2K_1$ 减小，总反应速率反而减小。这说明“升温速率一定增大”只对基元反应成立，对复杂反应要具体分析。
\end{jiexilist}
